\documentclass[11pt]{article}
\usepackage{amsmath,amssymb,amsfonts}
\usepackage[margin=1in]{geometry}
\usepackage{graphicx}
\usepackage{bm}
\usepackage{hyperref}

\title{Aeroelastic Vibration Analysis of a Finite-Span NACA~2412 Wing\\[4pt]
\large Version 2: FEM, Quasi-Steady Strip Theory, and Exploratory Unsteady Extensions}
\author{}
\date{}

\begin{document}
\maketitle

\section*{Version 2 Scope and Validity}

Version~2 is an extensible aeroelastic analysis framework rather than a validated
production flutter-prediction tool. The structural finite-element model and the
quasi-steady strip-theory branch provide a coherent qualitative and engineering
sensitivity model. The separate Theodorsen branch adds frequency-dependent circulatory
aerodynamics, a low-order added-mass approximation, a rational Theodorsen approximation,
and a Sears-based harmonic gust response.

The quasi-steady branch can be used to study trends in modal frequency, growth rate,
static trim, divergence indicators, and forced response. Absolute flutter speeds and
displacement amplitudes should not be interpreted as experimental predictions until the
structural properties, damping, aerodynamic data, and sign conventions have been
validated against a classical typical-section solution or measured NACA~2412 data.

The Theodorsen branch is exploratory. Its frequency-dependent matrices are evaluated
within a fixed-point $p$--$k$-style iteration, and its low-order added-mass and Sears
terms are approximations. Candidate eigenvalue real-part plots are diagnostic and are
not equivalent to a rigorously continued global nonlinear spectrum. Quantitative use
requires a validated rational-function state-space realization, complete unsteady
finite-wing aerodynamics, and benchmark comparison.

\section{Problem Setup}

The wing is idealized as a straight cantilever beam clamped at the root ($y=0$) and
free at the tip ($y=L$), where $L$ is the semi-span. Each spanwise station has two
elastic degrees of freedom: a plunge (bending) deflection $h(y,t)$, positive downward,
measured at the elastic axis, and a pitch (twist) rotation $\theta(y,t)$, positive
nose-up, about the same elastic axis. The chord distribution is $c(y)$ and the section
is a NACA~2412 profile at every station (no spanwise twist or taper assumed, though the
formulation supports both). The wing flies at freestream speed $U$ and air density
$\rho$, with dynamic pressure $q_\infty = \tfrac12 \rho U^2$.

The objective is to determine how the wing's vertical (plunge-dominated) vibration
behaves as $U$ is varied, for a wing of finite span $L$ (finite aspect ratio
$AR = b^2/S$, where $b=2L$ is the full span and $S$ the planform area).

\section{Structural Model}

\subsection{Continuous governing equations}

Euler--Bernoulli bending and Saint-Venant torsion of a non-uniform beam, each driven by
the distributed aerodynamic load:
\begin{align}
\frac{\partial^2}{\partial y^2}\!\left(EI(y)\,\frac{\partial^2 h}{\partial y^2}\right)
+ m(y)\,\frac{\partial^2 h}{\partial t^2}
&= L'(y,t), \label{eq:bending}\\[4pt]
-\frac{\partial}{\partial y}\!\left(GJ(y)\,\frac{\partial \theta}{\partial y}\right)
+ I_\theta(y)\,\frac{\partial^2 \theta}{\partial t^2}
&= M'_{ea}(y,t), \label{eq:torsion}
\end{align}
where $EI(y)$ is bending stiffness, $GJ(y)$ torsional stiffness, $m(y)$ mass per unit
length, $I_\theta(y)$ mass moment of inertia per unit length about the elastic axis,
$L'(y,t)$ the aerodynamic lift per unit span, and $M'_{ea}(y,t)$ the aerodynamic pitching
moment per unit span about the elastic axis.

\subsection{Finite element discretization}

The span is divided into $N_e$ two-node beam elements of length $L_e$. Each node
carries three degrees of freedom: $\{w,\ w',\ \theta\}$ (transverse deflection, slope,
twist). For a uniform element the classical consistent mass/stiffness sub-matrices are

\begin{equation}
\bm{K}^{(e)}_{\text{bend}} = \frac{EI}{L_e^3}
\begin{bmatrix}
12 & 6L_e & -12 & 6L_e\\
6L_e & 4L_e^2 & -6L_e & 2L_e^2\\
-12 & -6L_e & 12 & -6L_e\\
6L_e & 2L_e^2 & -6L_e & 4L_e^2
\end{bmatrix}, \qquad
\bm{M}^{(e)}_{\text{bend}} = \frac{mL_e}{420}
\begin{bmatrix}
156 & 22L_e & 54 & -13L_e\\
22L_e & 4L_e^2 & 13L_e & -3L_e^2\\
54 & 13L_e & 156 & -22L_e\\
-13L_e & -3L_e^2 & -22L_e & 4L_e^2
\end{bmatrix}
\end{equation}

acting on $\{w_1,w_1',w_2,w_2'\}$, and the rod-type torsion sub-matrices

\begin{equation}
\bm{K}^{(e)}_{\text{tors}} = \frac{GJ}{L_e}
\begin{bmatrix} 1 & -1\\ -1 & 1 \end{bmatrix}, \qquad
\bm{M}^{(e)}_{\text{tors}} = \frac{I_\theta L_e}{6}
\begin{bmatrix} 2 & 1\\ 1 & 2 \end{bmatrix}
\end{equation}

acting on $\{\theta_1,\theta_2\}$. Bending and torsion are structurally uncoupled
(elastic axis = reference axis), so the two blocks are merged into a single $6\times6$
element matrix ordered $\{w_1,w_1',\theta_1,w_2,w_2',\theta_2\}$ and assembled into
global matrices $\bm{M}_s,\bm{K}_s$ by standard FEM assembly. The three root DOFs are
removed (cantilever condition $w=w'=\theta=0$ at $y=0$), leaving a reduced DOF vector
$\bm{q}(t)$ of length $3N_e$ (minus 3 removed root DOFs $+3$ trailing-edge-of-root node
accounting -- i.e. $3N_e$ free DOFs after elimination).

\section{Aerodynamic Model: Quasi-Steady Strip Theory}

\subsection{Local effective incidence}

At spanwise station $y$, the local plunge velocity $\dot h(y,t)$ adds an induced angle
of attack, so the effective incidence seen by the strip is
\begin{equation}
\alpha_{\text{eff}}(y,t) = \theta(y,t) + \frac{\dot h(y,t)}{U}.
\label{eq:alphaeff}
\end{equation}
This is the baseline model implemented below and in the accompanying code. A
complete quasi-steady treatment also includes a pitch-\emph{rate}-dependent
circulatory lift term (Theodorsen's quasi-steady limit adds
$\left(\tfrac12-a\right)\tfrac{c(y)}{U}\dot\theta(y,t)$ inside
$\alpha_{\text{eff}}$, where $a$ locates the elastic axis as a fraction of
chord from mid-chord). That term's contribution to the lift is physically
distinct from the angle-of-attack term above, and critically, it does \emph{not}
act through the same moment arm $e(y)$ used in
Eq.~\eqref{eq:moment} below -- its own moment arm about the elastic axis
requires separate, careful derivation (see Theodorsen~1935, or Hodges \&
Pierce, \emph{Introduction to Structural Dynamics and Aeroelasticity}). Rather
than include that term on unverified footing, it is deliberately omitted from
the baseline formulation and code; Section~7 notes the resulting limitation
and a safe way to work around it.

\subsection{Section lift and moment}

\begin{align}
L'(y,t) &= q_\infty\, c(y)\, C_{L\alpha}\,\alpha_{\text{eff}}(y,t), \label{eq:lift}\\
M'_{ea}(y,t) &= e(y)\,L'(y,t) \;+\; q_\infty\, c(y)^2\, C_{m,ac}, \label{eq:moment}
\end{align}
where $e(y)$ is the offset of the aerodynamic center (quarter-chord) ahead of the
elastic axis, $C_{L\alpha}$ the lift-curve slope, and $C_{m,ac}$ the (camber-induced,
load-independent) moment coefficient about the aerodynamic center. The $C_{m,ac}$ term
is a constant, non-linear-in-$(q,\dot q)$ static twisting moment; it contributes to the
mean trim twist but \emph{not} to the linear dynamic stability matrices below, and is
tracked separately.

\subsection{Finite-span correction}

Equation~\eqref{eq:lift} uses the 2D (infinite-span) lift-curve slope. For a finite wing
of aspect ratio $AR$ and span efficiency $e_s$ (Oswald efficiency factor), the
effective 3D lift-curve slope from Prandtl lifting-line theory is
\begin{equation}
C_{L\alpha}^{3D} = \frac{C_{L\alpha}^{2D}}{1 + \dfrac{C_{L\alpha}^{2D}}{\pi\, AR\, e_s}}.
\label{eq:prandtl}
\end{equation}
This reduced slope is used uniformly in Eq.~\eqref{eq:lift} as a first-order finite-span
correction; a spanwise-varying correction (e.g. via a vortex-lattice induced-angle
solve) is a natural refinement.

\section{Aeroelastic Coupling: Generalized Equations of Motion}

The FEM shape functions interpolate $h(y,t)$ and $\theta(y,t)$ from the nodal DOF vector
$\bm{q}(t)$: $h(y,t) = \bm{N}_w(y)\,\bm{q}(t)$, $\theta(y,t) = \bm{N}_\theta(y)\,\bm{q}(t)$.

\textbf{Sign convention.} $h$ (and the nodal DOF $w$) is positive \emph{downward}
(Section~1), while $L'$ is the physical lift, positive \emph{upward} as is
standard in aerodynamics. The generalized force conjugate to $w$ is therefore
$-L'$: an upward lift resists an increase in downward deflection. This is the
classical typical-section convention (e.g.\ Hodges \& Pierce; Wright \&
Cooper), $m\ddot h + K h = -L'$, together with $I_\theta \ddot\theta + K\theta =
+M'_{ea}$ for the moment equation ($\theta$ positive nose-up, no sign flip
needed there).

The spanwise-distributed aerodynamic loads of
Eqs.~\eqref{eq:lift}--\eqref{eq:moment} are linear in $\bm{q}$ and $\dot{\bm q}$, so the
generalized aerodynamic force vector, lumped node-by-node (one aerodynamic
``strip'' per node, matching the discretization request this document
responds to) rather than integrated against the full FEM shape functions, can
be written as
\begin{equation}
\bm{Q}_{\text{aero}}(t) = -\,\bm{D}_a(U)\,\dot{\bm q}(t) \;-\; \bm{S}_a(U)\,\bm{q}(t)
\;+\; \bm{Q}_{\text{static}},
\end{equation}
with, at node $i$ (local block in the $\{w_i,w_i',\theta_i\}$ ordering, all
other entries zero),
\begin{equation}
\bm S_a^{(i)} = \begin{bmatrix} 0&0&+L'_{,\theta}\\ 0&0&0\\ 0&0&-M'_{,\theta}\end{bmatrix}\Delta y_i,
\qquad
\bm D_a^{(i)} = \begin{bmatrix} +L'_{,\dot h}&0&0\\ 0&0&0\\ -M'_{,\dot h}&0&0\end{bmatrix}\Delta y_i,
\end{equation}
where $L'_{,\theta}=q_\infty c\,C_{L\alpha}^{3D}$, $L'_{,\dot h}=q_\infty c\,
C_{L\alpha}^{3D}/U$, and $M'_{,(\cdot)} = e(y)\,L'_{,(\cdot)}$ (from
Eq.~\eqref{eq:moment}, dynamic part only), and $\Delta y_i$ is node $i$'s
tributary span width (half of each adjacent element). $\bm D_a$ is the
\emph{aerodynamic damping matrix} and $\bm S_a$ the \emph{aerodynamic stiffness
matrix}; both scale with $q_\infty \propto U^2$ except for the $\dot h/U$-based
entries, which scale linearly in $U$. The slope DOF $w'$ carries no direct
aerodynamic generalized force in this lumped, strip-per-node treatment -- a
consistent shape-function-weighted integral would add one; this is a
deliberate simplification matching the ``divide the span into sections''
formulation requested.

The coupled linear equations of motion are then
\begin{equation}
\bm M_s\,\ddot{\bm q} + \big(\bm C_s + \bm D_a(U)\big)\,\dot{\bm q}
+ \big(\bm K_s + \bm S_a(U)\big)\,\bm q = \bm 0,
\label{eq:coupled}
\end{equation}
where $\bm C_s$ is structural damping (often a small Rayleigh damping matrix,
$\bm C_s = \alpha \bm M_s + \beta \bm K_s$).

\section{Flutter / Stability Analysis (p-method)}

Assume a solution of the form $\bm q(t) = \hat{\bm q}\, e^{st}$. Substitution into
Eq.~\eqref{eq:coupled} gives the quadratic eigenvalue problem
\begin{equation}
\Big[s^2 \bm M_s + s\big(\bm C_s + \bm D_a(U)\big) + \big(\bm K_s + \bm S_a(U)\big)\Big]
\hat{\bm q} = \bm 0,
\label{eq:eig}
\end{equation}
solved at each trial airspeed $U$ (e.g. via linearization to a first-order state-space
form and a standard eigenvalue solver). Each eigenvalue $s_k(U) = \sigma_k(U) \pm
i\,\omega_k(U)$ gives, at that speed, a modal frequency $f_k = \omega_k/2\pi$ and a
damping ratio $g_k = -\sigma_k/\omega_k$ (or equivalently $\sigma_k$ itself).

\begin{itemize}
\item \textbf{V--f diagram}: $\omega_k(U)/2\pi$ vs.\ $U$ for each mode; flutter is
indicated where two modal frequencies approach one another (coalescence).
\item \textbf{V--g diagram}: $\sigma_k(U)$ (or $g_k$) vs.\ $U$; the \emph{flutter speed}
$U_F$ is the lowest $U$ at which any $\sigma_k$ crosses zero from negative (damped) to
positive (unstable).
\item \textbf{Torsional divergence} is a static ($\omega=0$) instability and shows up
separately as $\det(\bm K_s + \bm S_a(U)) \to 0$.
\end{itemize}

In the implementation, the complete state-space spectrum is retained by default rather
than truncating the search to a prescribed number of modes. Oscillatory branches are
continued between adjacent airspeeds using eigenvector modal assurance criterion (MAC)
correlation. If a branch becomes real or overdamped, its oscillatory frequency and growth
rate are recorded as unavailable ($\mathrm{NaN}$) rather than being replaced by another
mode. The real eigenvalues are retained and plotted separately. Dynamic instability is
reported directly from
\begin{equation}
 \sigma_{\max}(U)=\max_k \Re(s_k(U)),
\end{equation}
using the complete spectrum, including real eigenvalues; the first speed where
$\sigma_{\max}\geq 0$ is the reported instability speed.
Static divergence is checked independently using the normalized indicator
\begin{equation}
 d_D(U) =
 \frac{\det(\bm K_{s,r}+\bm S_{a,r}(U))}
 {\det(\bm K_{s,r})},
\end{equation}
where the subscript $r$ denotes the free (cantilever) DOFs. A sign change in $d_D$
brackets a divergence speed; the determinant indicator is not used as a substitute for
the dynamic flutter test.

For $U < U_F$, forced response to gusts/turbulence (replacing the RHS of
Eq.~\eqref{eq:coupled} with a stochastic or deterministic gust-force vector) gives the
finite vertical vibration amplitude at the tip, root, or any station of interest; this
amplitude grows sharply as $U \to U_F$ because the effective damping of the governing
mode approaches zero.

The companion code also provides a prescribed vertical-gust response. For a gust velocity
$w_g(t)$, the additional local incidence is $\alpha_g=w_g/U$, giving
\begin{equation}
 L'_g(y,t)=q_\infty c(y)C_{L\alpha}^{3D}\frac{w_g(t)}{U},\qquad
 Q_{w,g}=-L'_g\Delta y,\qquad Q_{\theta,g}=e(y)L'_g\Delta y.
\end{equation}
These generalized forces are assembled node-by-node and applied on the right-hand side
of the time-domain equations. The supplied example uses a sinusoidal gust; the same
interface accepts an impulse, one-minus-cosine gust, or stochastic gust function.

\section{NACA 2412 Aerodynamic Data}

\begin{itemize}
\item Thin-airfoil lift-curve slope: $C_{L\alpha}^{2D} = 2\pi\ \text{rad}^{-1}$
(theoretical); experimentally, finite-thickness/viscous effects typically give
$C_{L\alpha}^{2D} \approx 5.7$--$6.3\ \text{rad}^{-1}$ at low Mach, pre-stall.
\item Zero-lift angle of attack (due to camber): $\alpha_0 \approx -2.1^\circ$.
\item Moment coefficient about the aerodynamic center (quarter-chord), essentially
constant pre-stall: $C_{m,ac} \approx -0.045$ to $-0.05$.
\item Aerodynamic center location: $\approx 0.24$--$0.25\,c$ from the leading edge.
\end{itemize}
These feed directly into Eqs.~\eqref{eq:lift}--\eqref{eq:moment} as the per-section
aerodynamic properties, uniform along the span if the wing is untapered and untwisted.

\section{Numerical Implementation Correspondence}

The companion MATLAB code mirrors this document one section at a time:
\begin{enumerate}
\item \texttt{naca2412\_data.m} --- Section~6 aerodynamic coefficients.
\item \texttt{beam\_fem\_matrices.m} --- Section~2.2, assembling $\bm M_s,\bm K_s$.
\item \texttt{aero\_strip\_matrices.m} --- Section~3--4, building $\bm D_a(U)$,
$\bm S_a(U)$ from strip integration with the finite-span correction of
Eq.~\eqref{eq:prandtl}.
\item \texttt{flutter\_sweep.m} --- Section~5, sweeping $U$, solving
Eq.~\eqref{eq:eig}, retaining the complete spectrum, MAC-tracking oscillatory modes,
separating real eigenvalues, checking the static divergence indicator, and returning
the dominant eigenvalue.
\item \texttt{gust\_response.m} --- prescribed vertical-gust forcing and time-domain
integration of the coupled system below a selected flight speed.
\end{enumerate}

The parameter $a_{\rm ea}$ is passed through the aerodynamic interface but is unused by
the baseline model. It would enter the omitted pitch-rate circulatory contribution
$({1}/{2}-a_{\rm ea})c\dot\theta/U$. That contribution has a distinct aerodynamic
moment arm and is therefore not inserted using the angle-of-attack moment arm $e(y)$
without a separate checked derivation.

\section{MATLAB Code Block Guide}

This section describes the role of the principal code blocks in the MATLAB files and
their correspondence with the equations above.

\subsection{\texttt{naca2412\_data.m}: section data}

The function \texttt{naca2412\_data} creates the \texttt{aero} structure used by the
aerodynamic routines. The fields are:
\begin{itemize}
\item \texttt{Cl\_alpha\_2D}: the two-dimensional lift-curve slope;
\item \texttt{alpha0}: the measured zero-lift angle of attack;
\item \texttt{Cm\_ac}: the quarter-chord pitching-moment coefficient;
\item \texttt{ac\_loc}: the aerodynamic-center chordwise location.
\end{itemize}
The current linear stability matrices use \texttt{Cl\_alpha\_2D}, \texttt{Cm\_ac}, and
\texttt{ac\_loc}. The \texttt{alpha0} value is stored for static trim and extended
section-load calculations but is not included in the present dynamic matrices.

\subsection{\texttt{beam\_fem\_matrices.m}: structural assembly}

The input block defines the span mesh and evaluates any spanwise-varying structural
properties at each element midpoint. For every two-node element, the code then:
\begin{enumerate}
\item constructs the $4\times4$ Euler--Bernoulli bending stiffness and consistent mass
matrices;
\item constructs the $2\times2$ rod-type torsional stiffness and mass matrices;
\item embeds these submatrices in a $6\times6$ element matrix ordered as
 $\{w_1,w'_1,\theta_1,w_2,w'_2,\theta_2\}$;
\item scatters the element matrices into the global $\bm K_s$ and $\bm M_s$ matrices.
\end{enumerate}
The final block removes the three root DOFs and returns \texttt{free\_dofs}, which is
used to form the cantilever system for all subsequent calculations.

\subsection{\texttt{aero\_strip\_matrices.m}: strip loads}

This function first evaluates
\begin{equation}
 C_{L\alpha}^{3D}
 =
 \frac{C_{L\alpha}^{2D}}
 {1+C_{L\alpha}^{2D}/(\pi AR e_s)}
\end{equation}
and computes a tributary span width for every node. The loop over nodes evaluates the
local chord and the aerodynamic-center-to-elastic-axis offset
$e=(x_{ea}-x_{ac})c$. It then forms the lift sensitivities
\begin{equation}
 \frac{\partial L}{\partial\theta}
 =q_\infty c C_{L\alpha}^{3D}\Delta y,\qquad
 \frac{\partial L}{\partial\dot w}
 =\frac{q_\infty c C_{L\alpha}^{3D}}{U}\Delta y.
\end{equation}
These sensitivities are placed into the local $3\times3$ aerodynamic stiffness and
damping blocks. The slope DOF $w'$ receives no direct load in this lumped strip
approximation. The function also returns \texttt{Qstatic}, which contains the static
camber-induced lift and pitching-moment contributions and is separate from the linear
stability matrices.

\subsection{\texttt{flutter\_sweep.m}: eigenvalue sweep and tracking}

For each airspeed, the function:
\begin{enumerate}
\item reduces $\bm M_s$, $\bm K_s$, and $\bm C_s$ to the free DOFs;
\item obtains $\bm D_a(U)$ and $\bm S_a(U)$ from
 \texttt{aero\_strip\_matrices};
\item forms the state matrix
\begin{equation}
 \bm A(U)=
 \begin{bmatrix}
 \bm 0 & \bm I\\
 -\bm M_r^{-1}(\bm K_r+\bm S_{a,r}) &
 -\bm M_r^{-1}(\bm C_r+\bm D_{a,r})
 \end{bmatrix};
\end{equation}
\item computes the complete eigenvalue set $s_k(U)$;
\item stores the eigenvalue with the largest real part as the dominant stability
 eigenvalue;
\item separates real/overdamped eigenvalues from oscillatory eigenvalues;
\item uses eigenvector MAC correlation to continue oscillatory branches between adjacent
 airspeeds, leaving an unmatched branch as \texttt{NaN};
\item evaluates the normalized static-divergence determinant indicator.
\end{enumerate}
The returned \texttt{max\_growth} is $\max_k\Re(s_k)$ over the complete spectrum and
is therefore the quantity used for the first-instability check.

\subsection{\texttt{run\_flutter\_analysis.m}: top-level workflow}

The script is organized into five stages:
\begin{enumerate}
\item define geometry, structural properties, damping, and the airspeed grid;
\item construct the NACA data and structural matrices;
\item run the aeroelastic eigenvalue sweep;
\item report the first complete-spectrum instability and the separate divergence
 diagnostic;
\item plot the V--f, V--g, real-eigenvalue, and static-divergence diagrams.
\end{enumerate}
The final blocks select a check speed, integrate an initial-condition response using
\texttt{ode45}, and then call \texttt{gust\_response} for a prescribed sinusoidal
vertical gust.

\subsection{\texttt{gust\_response.m}: forced response}

The gust function converts a prescribed vertical gust velocity $w_g(t)$ into an
additional incidence $\alpha_g=w_g/U$. At every node it forms the gust lift
$L'_g=q_\infty c C_{L\alpha}^{3D}\alpha_g\Delta y$ and assembles the generalized
force components
\begin{equation}
 Q_{w,g}=-L'_g,\qquad Q_{\theta,g}=eL'_g.
\end{equation}
The reduced forced state equation is then integrated with \texttt{ode45}. The returned
\texttt{F\_hist} records the applied generalized gust force at the solver output times.

\textbf{Note on validation.} This is a first-principles quasi-steady strip-theory
template, not a pre-validated production tool. The matrix assembly was checked with an
independent re-implementation (structural $\bm M_s,\bm K_s$ symmetric and positive
definite on the free DOFs; the lowest cantilever bending/torsion frequencies match
beam-theory expectations; the plunge-dominated mode is stably damped at low airspeed
and becomes unstable only above a finite crossing speed, matching the expected
single-mode-becomes-unstable flutter pattern). Known limitations to be aware of before
trusting absolute numbers:
\begin{itemize}
\item \textbf{No torsional aerodynamic damping term} is included (Section~3.1):
the torsion-dominated mode can show a very weak instability even at low airspeed if
structural damping $\bm C_s$ is left at exactly zero. \texttt{run\_flutter\_analysis.m}
defaults to a small nonzero Rayleigh $\beta$ specifically to avoid this; replace it with
measured/estimated structural damping, or implement the omitted pitch-rate circulatory
term from a checked reference, for quantitative work.
\item No unsteady lag (Theodorsen's $C(k)$) or finite-span-varying induced angle
(vortex-lattice) is included -- both are noted as natural refinements in
Sections~3.3 and~7's companion code.
\item The gust response is a prescribed external-forcing calculation; it does not include
gust alleviation, unsteady aerodynamic lag, or a turbulence spectrum unless those are
supplied through the gust-velocity function.
\end{itemize}
Sanity-checking further against a classical closed-form case (e.g.\ Theodorsen's
flat-plate typical-section result) or a published NACA~2412 wind-tunnel/flutter test is
still recommended before using this for anything beyond a qualitative/teaching tool.

\end{document}
